% Generated by scripts/assemble.py; corrections belong in source/corrections.json.
\examhead{138}{华中师范大学 2024 年数学分析试卷}{131}
% p132-b20; source PDF page 132
\par\Needspace{5\baselineskip}
\blockanchor{p132-b20}
\majorquestion{一}计算题. 每题 10 分, 共 50 分.\par
% p132-b21; source PDF page 132
\par\Needspace{6\baselineskip}
\blockanchor{p132-b21}
\question{1}计算\par
\nopagebreak[4]
% p132-b22; source PDF page 132
\blockanchor{p132-b22}
\displayformula{\lim _ {x \rightarrow + \infty} \left[\left(x ^ {3} - x ^ {2} + \frac {x}{2}\right) e ^ {\frac {1}{x}} - \sqrt {1 + x ^ {6}} \right].}
% p132-b23; source PDF page 132
\blockanchor{p132-b23}
\question{2}将 $f(x) = \sin x, x \in [0, \pi]$ 展开成余弦级数.\par
% p132-b24; source PDF page 132
\blockanchor{p132-b24}
\question{3}计算 $\iint_{D} e^{\frac{y}{x + y}} \, \mathrm{d}x \, \mathrm{d}y$, 其中 $D = \{(x, y) \mid x + y \leq 1, x \geq 0, y \geq 0\}$.\par
% p132-b25; source PDF page 132
\par\Needspace{6\baselineskip}
\blockanchor{p132-b25}
\question{4}计算\par
\nopagebreak[4]
% p132-b26; source PDF page 132
\blockanchor{p132-b26}
\displayformula{I = \int_ {\Gamma} \frac {(x + y) \mathrm{d} x - (x - y) \mathrm{d} y}{x ^ {2} + y ^ {2}}.}
% p132-b27; source PDF page 132
\blockanchor{p132-b27}
其中  $\Gamma$  是由方程  $x^{2} + y^{2} = a^{2} (a > 0)$  确定的曲线, 方向取逆时针.\par
% p133-b00; source PDF page 133
\par\Needspace{6\baselineskip}
\blockanchor{p133-b00}
\question{5}计算\par
\nopagebreak[4]
% p133-b01; source PDF page 133
\blockanchor{p133-b01}
\displayformula{I = \iint_ {S} x \mathrm{d} x \mathrm{d} y + y \mathrm{d} z \mathrm{d} x + z \mathrm{d} x \mathrm{d} y.}
% p133-b02; source PDF page 133
\blockanchor{p133-b02}
其中 $S$ 是由方程 $x^{2} + y^{2} + z^{2} = a^{2}(a > 0)$ 确定的曲面，方向取外侧.\par
% p133-b03; source PDF page 133
\par\Needspace{5\baselineskip}
\par\Needspace{6\baselineskip}
\blockanchor{p133-b03}
\majorquestion{二}(15 分) 设 $\lim_{n\to \infty}a_n = a,\lim_{n\to \infty}b_n = b,$ 记\par
\nopagebreak[4]
% p133-b04; source PDF page 133
\blockanchor{p133-b04}
\displayformula{c _ {n} = \frac {a _ {1} b _ {n} + a _ {2} b _ {n - 1} + \cdots + a _ {n} b _ {1}}{n}.}
% p133-b05; source PDF page 133
\blockanchor{p133-b05}
请用数列极限的定义证明 $\lim_{n\to \infty}c_n = ab.$\par
% p133-b06; source PDF page 133
\par\Needspace{5\baselineskip}
\par\Needspace{6\baselineskip}
\blockanchor{p133-b06}
\majorquestion{三}(10 分) 设 $f(x)$ 在 $(- \infty, + \infty)$ 上二阶可微, 记 $M_{i} = \sup_{-\infty < x < + \infty}|f^{(i)}(x)| < + \infty (i = 0,1,2)$. 证明:\par
\nopagebreak[4]
% p133-b07; source PDF page 133
\blockanchor{p133-b07}
\displayformula{M _ {1} ^ {2} \leq 2 M _ {0} M _ {2}.}
% p133-b08; source PDF page 133
\par\Needspace{5\baselineskip}
\blockanchor{p133-b08}
\majorquestion{四}(15 分) 设 $(a, b)$ 为有界区间, 证明: 函数 $f(x)$ 在 $(a, b)$ 上一致连续的充要条件是对于 $(a, b)$ 上的任意 Cauchy 列 $\{x_{n}\}$, 函数值数列 $\{f(x_{n})\}$ 也是 Cauchy 列.\par
% p133-b09; source PDF page 133
\par\Needspace{5\baselineskip}
\blockanchor{p133-b09}
\majorquestion{五}(15 分) 设 $f(x)$ 为 $\mathbb{R}$ 上的连续函数, 记 $\varphi(x) = f(x) \int_{0}^{x} f(t) \, \mathrm{d}t, x \in (-\infty, +\infty)$, 若 $\varphi(x)$ 在 $(-\infty, +\infty)$ 上单调递减, 证明: $f(x) \equiv 0$.\par
% p133-b10; source PDF page 133
\par\Needspace{5\baselineskip}
\blockanchor{p133-b10}
\majorquestion{六}(10 分) 设 $\{x_{n}\}$ 为 $(0,1)$ 中的数列, 满足当 $i \neq j$ 时, $x_{i} \neq x_{j}$, 试讨论函数 $f(x) = \sum_{n=1}^{\infty} \frac{\operatorname{sgn}(x - x_{n})}{2^{n}}$ 在 $(0,1)$ 上的连续性, 其中 $\operatorname{sgn}(x)$ 表示符号函数.\par
% p133-b11; source PDF page 133
\par\Needspace{5\baselineskip}
\blockanchor{p133-b11}
\majorquestion{七}(15 分) 含参量反常积分 $I(r) = \int_{0}^{+\infty} r e^{-rx} \, \mathrm{d}x$.\par
% p133-b12; source PDF page 133
\blockanchor{p133-b12}
(1) 证明: $I(r)$ 在区间 $[a, b]$ 上关于 $r$ 一致收敛, 其中 $0 < a < b$.\par
% p133-b13; source PDF page 133
\blockanchor{p133-b13}
(2) 证明: $I(r)$ 在区间 $[0, b]$ 上关于 $r$ 不一致收敛, 其中 $0 < b$.\par
% p133-b14; source PDF page 133
\par\Needspace{5\baselineskip}
\blockanchor{p133-b14}
\majorquestion{八}（20分）设级数 $\sum_{n=1}^{\infty} a_n$ 收敛, 级数 $\sum_{n=1}^{\infty} (b_{n+1} - b_n)$ 绝对收敛.\par
% p133-b15; source PDF page 133
\blockanchor{p133-b15}
(1) 叙述阿贝尔变换, 即级数的分部求和公式.\par
% p133-b16; source PDF page 133
\blockanchor{p133-b16}
(2) 证明: $\lim_{n\to \infty}b_n$ 存在.\par
% p133-b17; source PDF page 133
\blockanchor{p133-b17}
(3) 证明: $\sum_{n=1}^{\infty} a_n b_n$ 收敛.\par
